% GENERATED by scripts/build-lecture-slides.py; do not edit. Portable profile: XeTeX, theme and fonts from the core.
\makeatletter\IfFileExists{core/beamer/beamerthememx.sty}{\def\mxcorepath{core/}\def\input@path{{core/beamer/}{core/}}}{\def\mxcorepath{../core/}\def\input@path{{../core/beamer/}{../core/}}}\makeatother
\documentclass[aspectratio=169,10pt,t]{beamer}
\usetheme{mx}
\usepackage{graphicx,booktabs}
\usepackage[french]{babel}
\setlength{\parskip}{.45em}
\hypersetup{pdfauthor={},pdfsubject={Lecture slides},pdfkeywords={},pdftitle={Algorithmique}}
\providecommand{\N}{\mathbb{N}}
\providecommand{\Z}{\mathbb{Z}}
\providecommand{\D}{\mathbb{D}}
\providecommand{\Q}{\mathbb{Q}}
\providecommand{\R}{\mathbb{R}}
\providecommand{\C}{\mathbb{C}}
\providecommand{\vect}[1]{\overrightarrow{#1}}
\providecommand{\vi}{\vec{\imath}}
\providecommand{\vj}{\vec{\jmath}}
\providecommand{\vk}{\vec{k}}
\providecommand{\norm}[1]{\left\lVert #1\right\rVert}
\providecommand{\vc}[2]{\begin{pmatrix}#1\\#2\end{pmatrix}}
\providecommand{\deter}[2]{\det\left(#1,#2\right)}
\providecommand{\abs}[1]{\left\lvert #1\right\rvert}
\providecommand{\dd}{\,\mathrm{d}}
\providecommand{\e}{\mathrm{e}}
\providecommand{\Card}{\operatorname{Card}}
\providecommand{\E}{\operatorname{E}}
\providecommand{\V}{\operatorname{V}}
\providecommand{\degre}{^{\circ}}
\providecommand{\un}[1]{\,\mathrm{#1}}
\providecommand{\sci}[2]{#1\times10^{#2}}
\providecommand{\pgcd}{\operatorname{PGCD}}
\providecommand{\ppcm}{\operatorname{PPCM}}
\title{Algorithmique}
\subtitle{Cours et exemples}
\date{}
\author{}
\begin{document}
\mxtitleframe
\begin{frame}{}
Un élève décrit oralement la construction d'une figure très simple sans la
montrer; un autre l'exécute. Comparer la figure obtenue à la figure cible et
repérer les instructions ambiguës. Réécrire la procédure avec des données
d'entrée, des instructions ordonnées et un résultat contrôlable.
\end{frame}
\begin{frame}{Objectifs}
\begin{itemize}

  \item Comprendre la notion d'algorithme et ses propriétés.
  
\item Lire, exécuter et tracer un algorithme en pseudo-code.
  
\item Tracer et lire un organigramme.
  
\item Écrire un algorithme simple pour un problème mathématique.

\end{itemize}
\end{frame}
\begin{frame}{Qu'est-ce qu'un algorithme ?}
Un \textbf{algorithme} est une suite finie et non ambiguë d'instructions
permettant de résoudre un problème. Il doit être :
\textbf{fini} (il termine), \textbf{valide} (il donne le bon résultat),
et \textbf{déterministe} (même entrée \(\Rightarrow\) même sortie).
\end{frame}
\begin{frame}{Variables, affectation, séquence}
Une \textbf{variable} est un emplacement mémoire identifié par un nom.
L'instruction \(A \leftarrow 5\) signifie   affecter la valeur \(5\) à \(A\) .
Attention : \(A \leftarrow A + 1\) remplace l'ancienne valeur de \(A\) par
\((\text{ancienne valeur}) + 1\).
\end{frame}
\begin{frame}{Variables, affectation, séquence}
Exécuter :
\(A \leftarrow 3,\quad B \leftarrow A + 7,\quad A \leftarrow A \times 2\)

\begin{center}\begin{tabular}{llll}
Étape & Instruction & \(A\) & \(B\) \\ 
1 & \(A \leftarrow 3\) & \(3\) & - \\ 
2 & \(B \leftarrow A+7\) & \(3\) & \(10\) \\ 
3 & \(A \leftarrow A \times 2\) & \(6\) & \(10\) \\ 
\end{tabular}\end{center}
\end{frame}
\begin{frame}{Structures de contrôle}
\begin{itemize}

  \item \textbf{Séquence} : les instructions s'exécutent dans l'ordre.
  
\item \textbf{Alternative} (\texttt{Si \ldots{} Sinon}) : exécution conditionnelle.
  
\item \textbf{Itération} (\texttt{Tant que} / \texttt{Pour}) : répétition.

\end{itemize}
\end{frame}
\begin{frame}{Méthodes et compléments de cours}
Pour \emph{exécuter à la main} un algorithme :
\end{frame}
\begin{frame}{Méthodes et compléments de cours (2)}
\begin{enumerate}
\setcounter{enumi}{0}
  \item Créer un tableau avec une colonne par variable.
  
\item Exécuter chaque instruction dans l'ordre, ligne par ligne.
  
\item À chaque affectation, mettre à jour la colonne concernée
        en barrant l'ancienne valeur.
  
\item Pour un \texttt{Si}/\texttt{Sinon}, évaluer la condition
        \emph{avant} de choisir la branche.
  
\item Pour un \texttt{Tant que}, répéter l'évaluation de la
        condition à chaque itération.

\end{enumerate}
\end{frame}
\begin{frame}{Méthodes et compléments de cours}
Dans ce manuel, on utilise les conventions suivantes :
\end{frame}
\begin{frame}{Méthodes et compléments de cours (2)}
\begin{itemize}

  \item \textbf{Affectation :} \(X \leftarrow \text{expression}\)
  
\item \textbf{Lecture :} \texttt{Lire}(\(X\))  (saisie utilisateur)
  
\item \textbf{Écriture :} \texttt{Afficher}(\(X\))
  
\item \textbf{Alternative :}
        \texttt{Si} \emph{condition} \texttt{alors} \ldots{}\textbackslash{}
        \texttt{Sinon} \ldots{} \texttt{Fin Si}
  
\item \textbf{Boucle bornée :}
        \texttt{Pour} \(i\) \texttt{de} \(a\) \texttt{à} \(b\) \texttt{faire}
        \ldots{} \texttt{Fin Pour}
  
\item \textbf{Boucle non bornée :}
        \texttt{Tant que} \emph{condition} \texttt{faire}
        \ldots{} \texttt{Fin Tant que}

\end{itemize}
\end{frame}
\begin{frame}{Méthodes et compléments de cours}
Un \textbf{invariant de boucle} est une propriété vraie avant
la boucle, conservée à chaque itération, et qui permet de
prouver la correction de l'algorithme à la sortie.
\end{frame}
\begin{frame}{Méthodes et compléments de cours}
Le \textbf{coût} d'un algorithme se mesure en comptant le nombre
d'opérations élémentaires effectuées en fonction de la taille \(n\)
de l'entrée. Plus ce nombre croît lentement avec \(n\), plus l'algorithme
est efficace. On parle d'algorithme \emph{rapide} (peu d'opérations) ou
\emph{lent} (beaucoup d'opérations).
\end{frame}
\end{document}
